
\pagebreak
\chapter{Logging.}


\medskip


Logging is an important component for the development and maintenance of the software applications.
It serves as the backbone for monitoring, diagnosing, and understanding application behavior in both development and production environments.

Logging is first of all important for debugging and diagnostics: logs provide insights into the execution flow and state of an application.
When an issue arises, logs act as breadcrumbs to help developers trace the root cause without needing to replicate the problem.

Logs are also important for performance monitoring, as by logging key metrics, such as response times and resource usage, developers can
identify performance bottlenecks and optimize application behavior.

Logs can detect and alert on suspicious activity, such as unauthorized access or anomalous behavior, contributing to the overall security of
the application.

Log levels categorize messages by their importance and urgency. Most systems use standard levels, such as:

\begin{itemize}
\item TRACE: Detailed information, typically of interest only during development or debugging.

\item DEBUG: Diagnostic information for debugging. Useful in non-production environments.

\item INFO: General informational messages about the application’s normal operations.

\item WARN: Indications of potential issues or non-critical problems that might need attention.

\item ERROR: Errors that prevent a specific operation from completing but do not crash the application.

\item FATAL: Severe errors that lead to application crashes or critical failures.
\end{itemize}

There are several best practices for introducing logging to your application.

First of all, we should awoid over-logging.
Logging everything can clutter log files and make it difficult to find useful information. Log only what’s necessary and at the appropriate level.

We should also consider contextual logging, meaning, that we should include relevant context in logs, such as user IDs, request IDs, or session IDs.
This makes logs more actionable and easier to correlate. But, we should never log sensitive information,
such as passwords or personal identifiable information (PII), unless properly masked or encrypted.

We should also regularly verify that log messages are meaningful, actionable, and formatted correctly.


After you have set up the logging in your application, good practice is to set up alerts, special monitoring tools to trigger alerts for crucial log
events, such as fatal errors or security violations.

In .net application logging can be implemented using either built-in tools, or NuGet packages, such as Serilog, NLog or log4net.

\medskip


\textbf{Implementation.} ASP.NET Core comes with a built-in logging framework (Microsoft.Extensions.Logging) that can log to console, debug, eventSource and third-party providers (Serilog, NLog, Application Insights, etc.)

When you create a new Web API project, logging is already configured in \verb"Program.cs":

\begin{lstlisting}
var builder = WebApplication.CreateBuilder(args);

// Add logging
builder.Logging.ClearProviders(); // optional: clear default providers
builder.Logging.AddConsole();     // log to console
builder.Logging.AddDebug();       // log to debug window

var app = builder.Build();
\end{lstlisting}

In order to add logging to the controller or service, you can inject \verb"ILogger<T>", where T is the service:


\begin{lstlisting}
[ApiController]
[Route("api/[controller]")]
public class ProductsController : ControllerBase
{
    private readonly ILogger<ProductsController> _logger;

    public ProductsController(ILogger<ProductsController> logger)
    {
        _logger = logger;
    }

    [HttpGet("{id}")]
    public IActionResult GetProduct(int id)
    {
        _logger.LogInformation("Fetching product with ID {ProductId}", id);

        if (id <= 0)
        {
            _logger.LogWarning("Invalid product id: {ProductId}", id);
            return BadRequest("Invalid ID");
        }

        try
        {
            // Simulate fetching product
            var product = new { Id = id, Name = "Test Product" };
            return Ok(product);
        }
        catch (Exception ex)
        {
            _logger.LogError(ex, "Error fetching product {ProductId}", id);
            return StatusCode(500, "Internal server error");
        }
    }
}

\end{lstlisting}


The built-in log levels can be used in the following way:

\begin{lstlisting}
_logger.LogTrace("Trace log");
_logger.LogDebug("Debug log");
_logger.LogInformation("Info log");
_logger.LogWarning("Warning log");
_logger.LogError("Error log");
_logger.LogCritical("Critical log");
\end{lstlisting}

You can control what gets logged without changing code, in the \verb"appsettings.json" file:

\begin{lstlisting}
{
  "Logging": {
    "LogLevel": {
      "Default": "Information",
      "Microsoft.AspNetCore": "Warning",
      "MyAppNamespace": "Debug"
    }
  }
}

\end{lstlisting}


For advanced logging (structured logs, file sinks, Seq, ElasticSearch), many teams use Serilog.

Install packages:

\begin{lstlisting}
dotnet add package Serilog.AspNetCore
dotnet add package Serilog.Sinks.Console
dotnet add package Serilog.Sinks.File

\end{lstlisting}

Configure in Program.cs:

\begin{lstlisting}
using Serilog;

var builder = WebApplication.CreateBuilder(args);

// Setup Serilog
Log.Logger = new LoggerConfiguration()
    .ReadFrom.Configuration(builder.Configuration) // read from appsettings.json
    .Enrich.FromLogContext()
    .WriteTo.Console()
    .WriteTo.File("logs/log-.txt", rollingInterval: RollingInterval.Day)
    .CreateLogger();

builder.Host.UseSerilog();

var app = builder.Build();

\end{lstlisting}

And in appsettings.json:


\begin{lstlisting}
"Serilog": {
  "MinimumLevel": "Information",
  "WriteTo": [
    { "Name": "Console" },
    { "Name": "File", "Args": { "path": "logs/webapi-.log", "rollingInterval": "Day" } }
  ]
}

\end{lstlisting}



